\documentclass[10pt]{article}  

%%%%%%%% PREÁMBULO %%%%%%%%%%%%
\title{Reporte de Laboratorio}
\usepackage[spanish]{babel} %Indica que escribiermos en español
\usepackage[utf8]{inputenc} %Indica qué codificación se está usando ISO-8859-1(latin1)  o utf8  
\usepackage{amsmath} % Comandos extras para matemáticas (cajas para ecuaciones,
% etc)
\usepackage{amssymb} % Simbolos matematicos (por lo tanto)
\usepackage{longtable} %agregadom para hacer tablas
\usepackage{xltabular}
\usepackage{graphicx} % Incluir imágenes en LaTeX
\usepackage{color} % Para colorear texto
\usepackage{subfigure} % subfiguras
\usepackage{float} %Podemos usar el especificador [H] en las figuras para que se
% queden donde queramos
\usepackage{capt-of} % Permite usar etiquetas fuera de elementos flotantes
% (etiquetas de figuras)
\usepackage{sidecap} % Para poner el texto de las imágenes al lado
	\sidecaptionvpos{figure}{c} % Para que el texto se alinie al centro vertical
\usepackage{caption} % Para poder quitar numeracion de figuras
\usepackage{commath} % funcionalidades extras para diferenciales, integrales,
% etc (\od, \dif, etc)
\usepackage{cancel} % para cancelar expresiones (\cancelto{0}{x})
 
\usepackage{anysize} 					% Para personalizar el ancho de  los márgenes
\marginsize{2cm}{2cm}{2cm}{2cm} % Izquierda, derecha, arriba, abajo

\usepackage{appendix}
\renewcommand{\appendixname}{Apéndices}
\renewcommand{\appendixtocname}{Apéndices}
\renewcommand{\appendixpagename}{Apéndices} 
% Para que las referencias sean hipervínculos a las figuras o ecuaciones y
% aparezcan en color
\usepackage[colorlinks=true,plainpages=true,citecolor=blue,linkcolor=blue]{hyperref}
%\usepackage{hyperref} 
% Para agregar encabezado y pie de página
\usepackage{fancyhdr} 
\pagestyle{fancy}
\fancyhf{}
\fancyhead[L]{\footnotesize UNI} %encabezado izquierda
\fancyhead[R]{\footnotesize Fac. de Ciencias}   % dereecha
\fancyfoot[R]{\footnotesize Reporte}  % Pie derecha
\fancyfoot[C]{\thepage}  % centro
\fancyfoot[L]{\footnotesize Reporte de conteo y medicion.}  %izquierda
\renewcommand{\footrulewidth}{0.4pt}


\usepackage{listings} % Para usar código fuente
\definecolor{dkgreen}{rgb}{0,0.6,0} % Definimos colores para usar en el código
\definecolor{gray}{rgb}{0.5,0.5,0.5} 
% configuración para el lenguaje que queramos utilizar
\lstset{language=Matlab,
   keywords={break,case,catch,continue,else,elseif,end,for,function,
      global,if,otherwise,persistent,return,switch,try,while},
   basicstyle=\ttfamily,
   keywordstyle=\color{blue},
   commentstyle=\color{red},
   stringstyle=\color{dkgreen},
   numbers=left,
   numberstyle=\tiny\color{gray},
   stepnumber=1,
   numbersep=10pt,
   backgroundcolor=\color{white},
   tabsize=4,
   showspaces=false,
   showstringspaces=false}

\newcommand{\sen}{\operatorname{\sen}}	% Definimos el comando \sen para el seno
%en español

\title{Plantilla para Reportes IMEC-UTB}
% Basada en la plantilla para reportes UPIITA de  Overleaf

%%%%%%%% TERMINA PREÁMBULO %%%%%%%%%%%%

\begin{document}

%%%%%%%%%%%%%%%%%%%%%%%%%%%%%%%%%% PORTADA %%%%%%%%%%%%%%%%%%%%%%%%%%%%%%%%%%%%%%%%%%%%
%%%
\begin{center}																		%%%
	\newcommand{\HRule}{\rule{\linewidth}{0.5mm}}									%%%\left
	%%%
	\begin{minipage}{0.48\textwidth} \begin{flushleft}
			\includegraphics[scale = 0.9]{Imagenes/UNI_2.png}
		\end{flushleft}\end{minipage}
	\begin{minipage}{0.48\textwidth} \begin{flushright}
			\includegraphics[scale = 0.075]{Imagenes/UNI.png}
		\end{flushright}\end{minipage}

	%%%
	\vspace*{1.0cm}								%%%
	%%%	
	\textsc{\huge Universidad Nacional de Ingieneria }\\[1.5cm]


	\begin{minipage}{0.9\textwidth}
		\begin{center}																					%%%
			\textsc{\LARGE Reporte  01 }
		\end{center}
	\end{minipage}\\[0.5cm]
	%%%
	%%%
	\vspace*{1cm}																		%%%
	%%%
	\HRule \\[0.4cm]																	%%%
	{ \huge \bfseries Laboratorio de Física I}\\[0.4cm]	%%%
	%%%
	\HRule \\[1.5cm]																	%%%
	%%%
	%%%
	\begin{minipage}{0.46\textwidth}													%%%
		\begin{flushleft} \large															%%%

			% Aqui a continuación pongan los nombres de los integrantes
			\emph{El grupo conformado por:}\\
			Rocca Cruz Axel Benjamin \\
			Llactahuaman Quispe Benjamin Israel \\
			Cortez Núñez Christian


			%%%
			%\vspace*{2cm}	
			%%%
			%%%
		\end{flushleft}																		%%%
	\end{minipage}
	%%%
	\begin{minipage}{0.52\textwidth}
		\vspace{-0.6cm}											%%%
		\begin{flushright} \large															%%%
			\emph{Profesores:} \\																	%%%
			Tello Gálvez Julio César\\
			Brocca Pobes Manuel Enrique%%%
		\end{flushright}																	%%%
	\end{minipage}
	\vspace*{1cm}
	%\begin{flushleft}

	%\end{flushleft}
	%%%
	\flushleft{\textbf{\Large Física BFI01}	}\\																		%%%
	\vspace{2cm}
	\begin{center}
		{\large \today}																	%%%
	\end{center}
\end{center}

\newpage
%%%%%%%%%%%%%%%%%%%% TERMINA PORTADA %%%%%%%%%%%%%%%%%%%%%%%%%%%%%%%%

\tableofcontents

\newpage

\begin{abstract}
	El desarrollo de este primer laboratorio de física está enfocado al estudio de los distintos procesos de medición y de guiar al estudiante a experimentar la inevitable aparición de errores o incertidumbres en cualquier medición y a aprender a tratar con ellos al momento de operar los resultados.
	La metodología de trabajo del laboratorio Nº 01 ha sido dividida en tres partes, siendo estas las siguientes:
	\begin{itemize}
		\item Medición y error experimental
		\item Medición y propagación de errores
		\item Gráfica de los resultados en una medición experimental
	\end{itemize}

\end{abstract}

\section{Bibliografía}
\begin{itemize}
	\item  Serway. F\'isica. Editorial McGraw-Hill (1992).\\
	\item Tipler. Física. Editorial Revert\'e (1994). \\
	\item Taylor, J.R. (2014) Introducci\'on al Análisis de errores - reverte, Taylor 2 Muestra.pdf. Disponible en: https://www.reverte.com/media/reverte/files/book-attachment-3746.pdf (Accessed: April 14, 2023).
\end{itemize}
\end{document}
